\documentclass[11pt]{article}
\usepackage[margin=1in]{geometry}
\usepackage{amsmath,amssymb,amsthm}
\usepackage[T1]{fontenc}
\usepackage{lmodern}
\usepackage[hidelinks]{hyperref}
\newtheorem{lemma}{Lemma}
\newtheorem{theorem}{Theorem}
\newcommand{\N}{\mathbb N}
\newcommand{\R}{\mathbb R}
\title{Conjecture 00000001741:\\disproof of the stated upward optimality}
\author{C0ldSmi1e}
\date{October 12, 2026}
\begin{document}
\maketitle
\begin{abstract}
The conjecture asserts infinitude of positive integer pairs satisfying a
cube--square gap bound, together with the assertion that its logarithmic
exponent cannot be increased. We prove that infinitude for this upper bound
is preserved by every increase of the exponent. Consequently the two stated
clauses cannot both hold. This does not decide the infinitude clause by itself,
Hall's conjecture, or a corrected optimality assertion involving a decrease
of the exponent.
\end{abstract}

\section{Exact assertion and formal interpretation}
The English conjecture reads:
\begin{quote}
The reverse direction of Hall's conjecture: there exist infinitely many pairs
with $0<|x^3-y^2|<\sqrt{x}\,(\log x)^{1/4}$; the logarithmic correction is optimal
(the exponent $1/4$ cannot be increased).
\end{quote}
The Chinese version also explicitly says that the exponent cannot be increased.
The exact bilingual source is included as \texttt{conjecture.md}. Both variables
are positive integers, as specified in the source's definition paragraph.
That paragraph provides background; no form of Hall's conjecture is assumed
as a hypothesis in this proof.

We use the natural logarithm and Mathlib's real powers. At boundary values,
these total functions use $\log 0=0$, $0^0=1$, and $0^a=0$ for $a\ne0$.
This makes the following definitions meaningful for every natural coordinate
and real exponent. The conjecture itself uses only positive coordinates and
positive exponents, where the usual nonnegative-base interpretation applies.
For $a\in\R$, write
\[
 B_a(x)=\sqrt{x}\,(\log x)^a,
 \qquad D(x,y)=|x^3-y^2|,
\]
\[
 S_a=\{(x,y)\in\N^2: x>0,\ y>0,\ 0<D(x,y)<B_a(x)\},
 \qquad I(a)\ \Longleftrightarrow\ S_a\text{ is infinite}.
\]
Subtraction in $D$ is over $\R$, after casting the natural-number coordinates;
it is not truncated subtraction. The condition $D>0$ excludes exact
cube--square coincidences. The statement is infinitude of distinct pairs.

The parenthetical specifies the direction of optimality. Its literal reading
in the displayed infinitude assertion gives the conjunction
\begin{equation}\label{eq:claim}
 C=I(1/4)\ \wedge\
       \bigl(\forall b\in\R,\ b>1/4\ \Longrightarrow\ \neg I(b)\bigr).
\end{equation}
We prove $\neg C$. The witness exponent will be $b=1$, so allowing only
rational larger exponents would not change the disproof. We do not replace
``increased'' by ``decreased'' or move the logarithm into a denominator.

\section{Finite exceptional pairs}
Increasing the exponent need not increase the bound when $0<\log x<1$.
Those small first coordinates must be treated separately, including their
possible second coordinates.

\begin{lemma}\label{lem:fiber}
For every real $a$ and fixed $x\in\N$, the set of pairs in $S_a$ with first
coordinate $x$ is finite. Consequently, for every $N\in\N$, the set of pairs
in $S_a$ with $x<N$ is finite.
\end{lemma}
\begin{proof}
Choose a natural number $M>x^3+B_a(x)$, using the Archimedean property of
$\R$. If $(x,y)\in S_a$, the strict absolute-value bound implies
\[
 y^2<x^3+B_a(x)<M.
\]
Since $y$ is a positive integer, $y\le y^2$, hence $y<M$. Thus this fiber is
contained in the finite set $\{(x,y):y\in\N,\ y<M\}$.
There are finitely many $x<N$, so the union of these finite fibers is finite.
\end{proof}

\begin{lemma}\label{lem:tail}
If $I(a)$ holds, then for every natural $N$ the set
$\{(x,y)\in S_a:N\le x\}$ is infinite.
\end{lemma}
\begin{proof}
Remove the finite set from Lemma~\ref{lem:fiber} from the infinite set $S_a$.
The remaining pairs have $x\ge N$.
\end{proof}

\section{Increasing the exponent preserves infinitude}
\begin{theorem}\label{thm:mono}
For all real $a\le b$, $I(a)$ implies $I(b)$.
\end{theorem}
\begin{proof}
Choose a natural $N\ge\exp(1)$. For $x\ge N$, monotonicity of the logarithm
gives $\log x\ge1$. Real exponentiation with base at least one is
nondecreasing in its exponent, so
\[
 (\log x)^a\le(\log x)^b,
 \qquad B_a(x)\le B_b(x).
\]
Multiplication by $\sqrt{x}\ge0$ preserves the inequality. Thus each pair in
$S_a$ with $x\ge N$ is also in $S_b$: the positivity and nonzero-gap
conditions remain true and $D(x,y)<B_a(x)\le B_b(x)$.
By Lemma~\ref{lem:tail}, that subset of $S_a$ is infinite if $I(a)$ holds.
Therefore $S_b$ is infinite.
\end{proof}

\begin{theorem}
The literal conjecture~\eqref{eq:claim} is false.
\end{theorem}
\begin{proof}
Assume $C$. Its first clause is $I(1/4)$. Theorem~\ref{thm:mono}, with
$a=1/4$ and $b=1$, gives $I(1)$. Since $1>1/4$, the second clause of $C$
gives $\neg I(1)$, a contradiction.
\end{proof}

This is a negation of a conjunction, not a construction of an infinite
approximation family and not a proof that $S_{1/4}$ is finite. In fact the
stronger implication $I(1/4)\Rightarrow I(b)$ is established for every
$b\ge1/4$. The choice of natural logarithm only fixes the finite cutoff;
the same argument works with any fixed logarithm base greater than one.

\clearpage
\section{Lean correspondence and reproduction}
The project uses Lean~4.19.0 and Mathlib~v4.19.0, with all transitive package
revisions pinned in \texttt{lean/lake-manifest.json}. The mathematical source
is \texttt{lean/HallReverse.lean}. In namespace \texttt{HallReverse},
\texttt{gap}, \texttt{upperBound}, \texttt{admissible}, and
\texttt{infinitelyMany} define $D$, $B$, $S$, and $I$ directly using
\texttt{Real.sqrt}, \texttt{Real.log}, real powers, and \texttt{Set.Infinite}.
The definitions \texttt{cannotIncrease} and \texttt{conjecture} give the
second clause and the full conjunction in~\eqref{eq:claim}.

The lemmas \texttt{finite\_fiber}, \texttt{finite\_below}, and
\texttt{infinite\_tail} prove the finite-exception argument.
The theorem \texttt{infinitelyMany\_mono} is Theorem~\ref{thm:mono}.
The final theorem is
\begin{center}
\texttt{HallReverse.conjecture\_false : $\neg$ HallReverse.conjecture}.
\end{center}
It has no extra assumptions. The source prints axiom dependencies for all
seven principal theorems.

\texttt{lean/Verification.lean} contains 13 kernel-checked examples: both
signs of the integer gap; exclusion of zero coordinates; exclusion of every
perfect cube--square coincidence; the zero bound and exclusion at $x=1$;
the fully expanded admissibility predicate; infinitude beyond any cutoff;
two specific larger exponents; general exponent monotonicity; and the
negation of the conjunction with every defining expression exposed.
Both modules are default Lake targets. From \texttt{lean/}, a standard
reproduction is:
\begin{verbatim}
lake exe cache get
lake --wfail build
lake env lean -DwarningAsError=true HallReverse.lean
lake env lean -DwarningAsError=true Verification.lean
\end{verbatim}

The package includes local verification records and an independent AI-agent
semantic review. These are contributor checks, not maintainer acceptance.
Verification covers a fresh build of the submitted modules, direct source
replay with warnings treated as errors, dependency-closure and axiom checks,
and explicit positive and negative controls for the auditing tool. The only
permitted foundational axioms are \texttt{propext}, \texttt{Classical.choice},
and \texttt{Quot.sound}; there are no admitted proofs, custom axioms, or
\texttt{native\_decide} proofs in the submitted mathematical modules.
Archived negative controls intentionally contain rejected constructs to test
the auditing tools; they are outside the mathematical proof library.
No auxiliary numerical computation is needed for the mathematics.

Pinned dependency sources and supplied cached imports are checked by identity;
the complete Mathlib library is not rebuilt, and the supplied Lean executable
is not bootstrapped. The source, tests, report, and recorded file hashes allow
the local validation to be distinguished from those trust assumptions.
\end{document}
